Chain-of-thought (CoT) prompting asks the model to produce intermediate reasoning steps before the final answer, which can substantially improve performance on arithmetic, commonsense, and symbolic reasoning tasks~\cite{wei2022chain}. Wei et al.\ showed that providing a few CoT exemplars in the prompt elicits step-by-step reasoning in sufficiently large models and can surpass direct-answer baselines on benchmarks such as GSM8K. We compare direct-answer and CoT prompting on a small set of arithmetic word problems using two contemporary models.

\paragraph{(a) Domain and problems.}
We chose arithmetic and short word problems and constructed 12 problems with known correct numeric answers covering multi-step arithmetic, discounts, rates, and geometry. The full list with answers is in Appendix~\ref{app:cot-data}.

\paragraph{(b) Direct-answer prompt.}
We wrote a prompt that asks the model to read the problem and give only the final numeric answer with no intermediate reasoning. The template is in Appendix~\ref{app:cot-prompts}.

\paragraph{(c) Chain-of-thought prompt.}
We wrote a second template that asks the model to show intermediate reasoning steps and then give the final numeric answer on the last line. The template is in Appendix~\ref{app:cot-prompts}.

\paragraph{(d) Evaluation of two models.}
We ran both prompts on all 12 problems with Gemini~2.5 Flash and Claude Haiku, recording correctness for each. Table~\ref{tab:cot-accuracy} reports accuracy along with macro precision, recall, and F1 for each model and prompting style.

\paragraph{(e) Comparison of prompting styles.}
For Gemini, direct-answer performed better (0.83) than CoT (0.75); CoT introduced arithmetic slips that propagated to the final answer. For Claude, both styles gave 0.75 accuracy, but CoT slightly improved macro precision and F1 as shown in Table~\ref{tab:cot-accuracy}. Chain-of-thought did not improve raw accuracy for either model and hurt one, though it can still improve the structure of errors without changing the overall score.

\begin{table}[h]
  \centering
  \begin{tabular}{lccccc}
    \toprule
    Model & Style & Accuracy & Precision & Recall & F1 \\
    \midrule
    Gemini 2.5 Flash & Direct & 0.83 & 0.64 & 0.64 & 0.64 \\
    Gemini 2.5 Flash & CoT    & 0.75 & 0.50 & 0.50 & 0.50 \\
    Claude Haiku     & Direct & 0.75 & 0.50 & 0.50 & 0.50 \\
    Claude Haiku     & CoT    & 0.75 & 0.58 & 0.56 & 0.57 \\
    \bottomrule
  \end{tabular}
  \caption{Direct-answer vs.\ chain-of-thought prompting (Problem~2).}
  \label{tab:cot-accuracy}
\end{table}